%%%%%%%%%%%%%%%%%%%%%%%%%%%%%%%%%%%%%%%%%%%%%%%%%%%%%%%%%%%%%
%% Examtask 2: SEPIC converter as smart voltage stabilizer for bicycle dynamo %%
%%%%%%%%%%%%%%%%%%%%%%%%%%%%%%%%%%%%%%%%%%%%%%%%%%%%%%%%%%%%%

\task{SEPIC converter as smart voltage stabilizer for bicycle dynamo}

\taskGerman{SEPIC Konverter als intelligente Spannungsstabilisierung für Fahrraddynamos}

A SEPIC converter is to be designed for voltage stabilization in a bicycle dynamo.
A control circuit regulates the converter's duty cycle to ensure a stable output voltage.
The dynamo provides a voltage of $\SI{3}{\volt} \ldots \SI{10}{\volt}$. The output voltage is to be stabilized at $\SI{5.1}{\volt}$.

\begin{germanblock}
    Für ein Fahraddynamo soll ein SEPIC-Konverter als Spannungsstabilisierung dimensioniert werden.
    Eine Reglungsschaltung regelt das Tastverhältnis des Konverter für eine stabile Ausgangsspannung.
    Der Dynamo liefert eine Spannung von $\SI{3}{\volt} \ldots \SI{10}{\volt}$. Die Ausgangsspannung soll auf $\SI{5.1}{\volt}$ stabilisiert werden.
\end{germanblock}

\input{fig/Task02/figSEPICConverter}

\begin{table}[ht]
    \centering  % Zentriert die Tabelle
    \begin{tabular}{llll}
        \toprule
        \multicolumn{2}{l}{\textbf{General parameters:}} \\ 
        Input voltage: &  $U_{\mathrm{1}} = \SI{3}{\volt} \ldots \SI{10}{\volt}$ & Diode forward voltage: & $U_{\mathrm{D},\mathrm{f}} = \SI{0.8}{\volt}$ \\
        Output voltage: & $U_{\mathrm{2}} = \SI{5}{\volt}$  & Output current range: & $I_2 = \SI{50}{\milli\ampere} \ldots \SI{1}{\ampere}$  \\
        Switching frequency: & $f_\mathrm{s} = \SI{100}{\kilo\hertz}$  &  &   \\
        \midrule
        \multicolumn{4}{l}{\quad The diode forward voltage must be taken into account.}  \\ 
        \multicolumn{4}{l}{\quad The voltage of capacitor $C_{\mathrm{1}}$ can be considered as constant for the specific duty cycles.}  \\ 
        \multicolumn{4}{l}{\quad The output voltage can be considered as constant except of subtask 2.4.}  \\ 
        \bottomrule
    \end{tabular}
    \caption{Parameters of the circuit.}  % Beschriftung der Tabelle
    \label{table:ex01_Parameters of the circuit}
\end{table}


\subtask{What are the minimum and maximum duty cycles that the control circuit has to request?}{0}
\subtaskGerman{Was ist das minimale und maximale Tastverhältnis, den die Reglungsschaltung vorgeben muss.}

\subtask{Which operation point is to consider, for the calculation of minimal $L_{\mathrm{1}}$, 
         to avoid DCM mode within the specified operating range? \\
         Determine the minimal required inductances $L{\mathrm{1}}$ for this case.}{0}
\subtaskGerman{Welcher Arbeitspunkt ist bei der Berechnung der minimalen Induktivität $L{\mathrm{1}}$ zu berücksichtigen, 
               um den DCM-Betrieb innerhalb des definierten Arbeitsbereichs zu vermeiden? \\
               Berechnen Sie die minimal erforderliche Induktivität $L{\mathrm{1}}$ für diesen Fall.}

\subtask{Which operation point is to consider, for the calculation of minimal $L_{\mathrm{2}}$, 
         to avoid DCM mode within the specified operating range? \\
         Calculate the minimal required inductances $L{\mathrm{2}}$ for this case.}{0}
\subtaskGerman{Welcher Arbeitspunkt ist bei der Berechnung der minimalen Induktivität $L{\mathrm{2}}$ zu berücksichtigen, 
               um den DCM-Betrieb innerhalb des definierten Arbeitsbereichs zu vermeiden? \\
               Berechnen Sie die minimal erforderliche Induktivität $L{\mathrm{1}}$ für diesen Fall.}

\subtask{How large must the filter capacitor be so that, in the worst case, the voltage fluctuation is $\pm \SI{5}{\percent}$ ?}{0}
\subtaskGerman{Wie groß muss der Filterkondensator ausgelegt werden, damit im ungünstigsten Fall die Spannungsschwankung $\pm \SI{5}{\percent}$ ?}
